\section{Introduction}
\label{sec:intro}

A common pattern in zk-friendly arithmetizations of bit-oriented
computations such as SHA-256 is the appearance, inside a single
constraint, of bit-wise rotations and right-shifts of one of the trace
columns. Two examples from SHA-256:
\begin{align*}
\sigma_0(x) &= \Rot_7(x) \oplus \Rot_{18}(x) \oplus \ShR_3(x), \\
\sigma_1(x) &= \Rot_{17}(x) \oplus \Rot_{19}(x) \oplus \ShR_{10}(x).
\end{align*}

\paragraph{The naive arithmetization.} Treat each rotated/shifted
operand as a fresh trace column $w := \bitop_*(v)$, commit to it
alongside $v$, and prove $w_i = \bitop(v_i)$ row-wise via additional
constraints. This costs:
\begin{itemize}[leftmargin=2em]
  \item one extra committed column per occurrence of $\bitop$;
  \item one extra equality constraint per occurrence;
  \item the corresponding extra opening, MLE evaluation, and proof bytes
    in every step downstream.
\end{itemize}
For SHA-256, where $\sigma_0$ and $\sigma_1$ are evaluated on every
message-schedule row and contribute six rotation/shift instances per
compression step, this overhead dominates the per-row witness budget
and obscures the algebraic structure of the underlying computation.

\paragraph{Our observation.} A binary-polynomial cell is an element of
$\cellring := \F_2[X]/(X^{32})$. The bit-wise operations
$\Rot_c$ and $\ShR_c$ act on cells as $\F_2$-linear endomorphisms, and
they act \emph{only on coefficient positions}: they do not depend on
the row index. The multilinear extension that lifts a column $v \in
\cellring^n$ to a polynomial $\mle{v}$ over the row hypercube is
$\F_2$-linear in the cell values. Because the cell coordinate and the
row coordinate are independent, $\bitop$ and the MLE-evaluation
operator commute:
\[
  \mle{\bitop_* v}(r) \;=\; \bitop\!\left(\mle{v}(r)\right)
  \qquad \text{for all } r \in \F_q^{\nu}.
\]

\paragraph{Implication.} A virtual column $w := \bitop_* v$ does not
need a commitment, an opening, or a fresh challenge. Its value at the
final evaluation point of the protocol is determined by the source
column's lifted opening, which the verifier already receives. The
prover sends a single field element per occurrence (the value at the
constraint-aggregation challenge $r^*$) so that the constraint
polynomial can be evaluated at $r^*$; soundness of that single field
element is bound to the source column through the existing
multi-point evaluation argument and the closed-form reconstruction at
the final point.

The construction is described in
\Cref{sec:setting,sec:equivariance,sec:construction}, its soundness
sketched in \Cref{sec:soundness}, its communication cost compared to
the naive arithmetization in \Cref{sec:cost}, and its $\F_2$-linear
generalisation in \Cref{sec:generalisation}.
